%===============================================================================
\section{Introduction}
\label{sec:introduction}
%===============================================================================

Understanding how cells respond to genetic perturbations, such as those induced by CRISPR screens~\citep{dixit2016perturb}, is fundamental to understand disease biology and advancing therapeutic discovery ~\citep{plenge2013validating,adamson2016multiplexed,norman2019exploring}. 
However, the combinatorial explosion across $\sim$25,000 genes makes exhaustive perturbation experiments infeasible~\citep{norman2019exploring}, motivating computational models that predict cellular responses to unseen perturbations or combinations thereof.

Existing approaches to predict the transcriptional outcome following a perturbation face a fundamental challenge: generative quality comes at the cost of causal interpretability.
For example, recent methods like GEARS~\citep{roohani2024predicting} or scGPT~\citep{cui2024scgpt} achieve strong predictive performance in cellular response to perturbations; however, they act as ``black boxes'' providing little mechanistic insight into how perturbations affect cellular systems,
%This opacity prevents researchers from understanding why certain combinations of perturbations produce synergistic effects, 
failing to provide testable biological hypotheses.


Computational biology demands models that provide mechanistic insights into cellular function~\citep{chen_applying_2024}, making \gls{crl}~\citep{scholkopf2021toward} a promising approach for mapping (intervened) genes to biological processes.
Recent \gls{crl} approaches such as \gls{dvae}~\citep{zhang2023identifiability} or SENA~\citep{de2025interpretable} aim at recovering the underlying causal factors aligned with the interventions, enabling counterfactual reasoning.
However, these \gls{crl} approaches suffer from a fundamental architectural limitation: a dense interventional encoder allows every gene, including the intervened ones, to influence many or all latent variables, undermining the requirements for their identifiability guarantees.
We term this phenomenon \emph{intervention spillover}, which renders identifying the unique latent target for a given intervention ill-posed.
Consequently, the intervention mechanism cannot perform the clean, single-node manipulation required for valid causal inference~\citep{brehmer2022weakly,lippe2022citris}.

\revised{Further, the destructive nature of \gls{scRNAseq} makes it impossible to observe the same cell before and after a genetic intervention, complicating the prediction of cell-specific responses.
Existing methods address this issue by randomly pairing control and perturbed cells.
Since current models are trained to minimize the error across all possible arbitrary pairings, they are incentivized to (suboptimally) learn the population-mean response.
We refer to this phenomenon as \emph{mean collapse}, as this regression to the mean erodes cellular diversity.}

To address these critical issues, we present \gls{raptorgraph}, a novel framework that addresses causal identifiability challenges through a structured encoder design. 
We first propose a novel GraphPathway layer, which uses preconditioning to enforce sparse mappings from genes to learned interpretable causal factors. We term these factors \emph{\cmps{}}, which encompass the (related) known biological processes associated with a given perturbation.
Our design enables clean, single-node latent associations required for causal inference while  maintaining interpretability.
This innovation is complemented by an \gls{ot} preprocessing step, which mitigates the distinct problem of mean collapse.

\revised{We benchmark \glsentryshort{raptorgraph} on two well-known datasets: the \citet{norman2019exploring} Perturb-seq dataset, and the \citet{replogle_combinatorial_2020} CRISPRi dataset, assessing, among others, the accuracy of inferred genetic interactions under combinations of perturbations, and the biological plausibility of the inferred \cmps{}. We uncover biologically accurate programs, such as differentiation-induced cell cycle arrest driven by EGR1 and the stress-induced growth arrest mediated by JUN.}
%\revised{Additionally, we \citet{replogle_combinatorial_2020} CRISPRi dataset}

%We demonstrate strong performance alongside interpretability, which addresses key gaps not filled by current approaches.
